\documentclass[11pt]{article}
\usepackage[T1]{fontenc}
\usepackage{lmodern}
\usepackage[margin=0.88in]{geometry}
\usepackage{amsmath,amssymb,amsthm,mathtools}
\usepackage{booktabs,longtable,array,tabularx}
\usepackage{xcolor,microtype,fancyhdr,enumitem}
\usepackage[hyphens]{url}
\usepackage{xurl}
\usepackage{listings}
\usepackage{graphicx}
\usepackage[hypertexnames=false,colorlinks=true,linkcolor=blue!45!black,
 citecolor=blue!45!black,urlcolor=blue!45!black]{hyperref}
\setlength{\headheight}{14pt}
\setlength{\emergencystretch}{3em}
\setlist{itemsep=3pt,topsep=5pt}
\lstset{basicstyle=\ttfamily\footnotesize,breaklines=true,
 columns=fullflexible,keepspaces=true,frame=single,rulecolor=\color{black!18}}
\pagestyle{fancy}\fancyhf{}
\fancyhead[L]{\small Lambert \(W\) in intuitionistic PRA}
\fancyhead[R]{\small constructive branches, series, and algorithms}
\fancyfoot[C]{\thepage}

\newtheorem{theorem}{Theorem}[section]
\newtheorem{lemma}[theorem]{Lemma}
\newtheorem{proposition}[theorem]{Proposition}
\newtheorem{corollary}[theorem]{Corollary}
\newtheorem{claim}[theorem]{Claim}
\newtheorem{example}[theorem]{Example}
\theoremstyle{definition}
\newtheorem{definition}[theorem]{Definition}
\newtheorem{remark}[theorem]{Remark}
\newtheorem{audit}[theorem]{Audit finding}
\newtheorem{algorithm}[theorem]{Algorithm}

\newcommand{\N}{\mathbb N}
\newcommand{\Z}{\mathbb Z}
\newcommand{\Q}{\mathbb Q}
\newcommand{\R}{\mathbb R}
\newcommand{\C}{\mathbb C}
\newcommand{\iPRA}{\mathsf{iPRA}}
\newcommand{\PR}{\mathsf{PR}}
\newcommand{\Wz}{W_0}
\newcommand{\Wm}{W_{-1}}
\newcommand{\Lambert}{W}
\newcommand{\abs}[1]{\left|#1\right|}
\newcommand{\norm}[1]{\left\lVert#1\right\rVert}
\newcommand{\apart}{\mathrel{\#}}
\newcommand{\code}[1]{\nolinkurl{#1}}
\newcommand{\coeff}[2]{[X^{#1}]\,#2}
\newcommand{\eps}{\varepsilon}
\newcommand{\Ei}{\operatorname{Ei}}
\newcommand{\sgn}{\operatorname{sgn}}
\newcommand{\cis}{\operatorname{cis}}
\newcommand{\StatusF}{\textsf{F}}
\newcommand{\StatusA}{\textsf{A}}
\newcommand{\StatusC}{\textsf{C}}
\newcommand{\StatusP}{\textsf{P}}
\newcommand{\StatusD}{\textsf{D}}

\title{\textbf{The Lambert \(W\) Function in Intuitionistic Primitive Recursive Arithmetic}\\
\large Real branches, finite Lagrange inversion, calculus, inequalities, asymptotics, algorithms, and applications}
\author{A constructive proof package prepared for Vladimir Reshetnikov}
\date{16 November 2026}
\hypersetup{pdftitle={The Lambert W Function in Intuitionistic Primitive Recursive Arithmetic},
 pdfauthor={Prepared for Vladimir Reshetnikov},
 pdfsubject={Constructive proofs and a Lean formalization roadmap for the Lambert W function}}

\begin{document}
\maketitle

\begin{abstract}
We reconstruct a large part of the mathematical content of the Wikipedia article
``Lambert W function'' in intuitionistic primitive recursive arithmetic
(\(\iPRA\)).  The phrase ``in \(\iPRA\)'' is made literal: real and complex
numbers are fixed primitive-recursive rapidly convergent rational names;
strict inequalities and division carry positive separation witnesses; analytic
functions are evaluation algorithms with explicit primitive-recursive moduli;
and every infinite assertion is translated into a first-order statement about
finite approximations.

The real branches are constructed by certified dyadic search for the inverse of
\(F(w)=w e^w\).  We prove their domains, ranges, uniqueness, monotonicity,
curvature, defining identities, special values, derivative and higher-derivative
formulae, elementary antiderivatives, and several definite integrals.  The
finite Lagrange--Buermann theorem from the companion ProveIt development gives
the exact Maclaurin coefficients, all integral powers, the analytic powers of
\(W_0(x)/x\), the radius \(1/e\), and the rooted-tree generating function.
Finite reversion at the critical point gives the square-root expansion.
Formal reversion of \(w+\log w=\log x\), followed by an explicit residual
estimate, gives every finite order of the logarithmic asymptotic expansion.

We give short constructive proofs of the Hoorfar--Hassani bounds, the sharper
Iacono--Boyd bounds, the square/cube-root bounds on \([-1/e,0]\), and both
families of bounds for \(W_{-1}\).  Newton, Halley, logarithmic Newton, bisection,
and interval-Newton evaluation are reduced to bounded rational algorithms with
certified error bounds.  Algebraic substitutions prove the standard formulas
for intersections of lines and exponentials, the inverse of \(x^x\), the
Bernoulli--Todd equation, orthogonal trajectories, and the general equation
\(e^{-cx}=a_0(x-r)\).

The archive contains this paper, a theorem ledger, exact symbolic and numerical
checks, and two detailed Lean blueprints tied to the current ProveIt
\code{ConciseIPRA} interfaces.  The Lean files are implementation outlines with
explicit \code{sorry} obligations, not a completed kernel-checked
formalization.  Claims depending on a global complex branch convention,
Lindemann--Weierstrass, specialized integral representations, or an external
physical model are separated from the native \(\iPRA\) results.
\end{abstract}

\newpage
\begingroup\small\tableofcontents\endgroup
\newpage

\section{Scope, status labels, and the main constructive theorem}\label{sec:scope}

The reference page mixes finite algebra, real inverse-function facts, complex
analytic continuation, inequalities, asymptotics, numerical recipes,
combinatorics, and statements about external scientific models.  These do not
all have the same logical status.  We use the following labels throughout.

\begin{center}
\begin{tabularx}{\textwidth}{>{\bfseries}c X}
\toprule
\StatusF & A finite identity or bounded algorithm, uniform in its natural-number
parameters, directly formalizable by quantifier-free induction in \(\iPRA\).\\
\StatusA & A theorem about fixed primitive-recursive real or complex names,
with explicit bounds and convergence moduli.  Substitution of the names gives a
first-order \(\iPRA\) sentence.\\
\StatusC & The reduction is constructive and arithmetized, but the conclusion
uses an additional named theorem not developed here (for example an effective
Lindemann--Weierstrass theorem).\\
\StatusP & A source-specific claim is reconstructed only partially or checked
computationally; the missing uniform certificate is stated explicitly.\\
\StatusD & Not claimed: a bibliographic, software, or physical-model statement,
or a specialized representation whose analytic certificate is not reconstructed
in this archive.\\
\bottomrule
\end{tabularx}
\end{center}

\begin{theorem}[Constructive Lambert atlas]\label{thm:atlas}
For every fixed collection of primitive-recursive input names and every positive
separation certificate required by a denominator or branch condition, \(\iPRA\)
proves the translations of the following statements.
\begin{enumerate}
\item The complete real-branch classification for \(w e^w=x\), including
constructive names for \(W_0\) and \(W_{-1}\), their domains, ranges,
monotonicity, continuity, and inverse identities.
\item The differential equation, the first and higher derivative formulae,
curvature, elementary antiderivatives, and the displayed definite integrals
whose auxiliary gamma or exponential-integral names are supplied with their
standard effective integral definitions.
\item The Maclaurin series, all finite coefficient identities for powers, the
exact radius \(1/e\), boundary convergence, the branch-point Puiseux series,
and every fixed finite order of the logarithmic asymptotic expansion.
\item The bounds in Sections~\ref{sec:bounds-principal} and
\ref{sec:bounds-secondary}, the algebraic identities and special values in
Sections~\ref{sec:identities} and \ref{sec:special}, and the equation-solving
and combinatorial applications in Sections~\ref{sec:equations} and
\ref{sec:trees}.
\item Certified bisection and interval-Newton evaluation on both real branches,
and Newton, Halley, and logarithmic Newton on any rational interval carrying a
positive lower bound for the relevant derivative.
\end{enumerate}
All uniformity is in numerical codes and bounded folds; there is no
quantification over an open-ended class of algorithms.
\end{theorem}

\begin{proof}[Proof architecture]
Sections~\ref{sec:meaning}--\ref{sec:realbranches} define the number codes,
construct \(e^x\) and \(\log x\), and invert the monotone pieces of
\(w\mapsto we^w\) by a primitive-recursive interval search.  The identities and
calculus then use finite algebra and explicit limit moduli.  The power-series
claims are consequences of a finite Lagrange--Buermann identity in
\(\Q[X]/(X^{N+1})\); the analytic claims add explicit majorants.  Each
inequality is reduced to monotonicity of \(F\) and one of a small set of
rationally certified logarithm inequalities.  The numerical algorithms are
bounded recursions whose invariants are interval inclusions.  Section
\ref{sec:arithmetization} spells out the object-theory translation.
\end{proof}

\subsection{The central audit boundary}

Three qualifications prevent overstatement.
\begin{enumerate}
\item A complex branch \(W_k\) is represented only after fixing a rational
base point, a root enclosure there, and a polygonal continuation path separated
from the critical value and the chosen cut.  This is enough for local branch
identities but does not silently import the classical Riemann-surface picture.
\item The page's formula \(W_0(-\pi/2)=i\pi/2\) lies on the conventional
principal cut.  Both \(i\pi/2\) and \(-i\pi/2\) solve \(we^w=-\pi/2\).  The
printed equality is therefore a boundary-value convention, not an
unqualified single-valued identity.
\item A displayed formula copied from a physical model is proved here only as
an algebraic solution of the displayed equation.  The claim that the equation
correctly models a fluid, circuit, epidemic, or quantum system requires the
external scientific derivation and is not an arithmetic consequence of the
formula.
\end{enumerate}

\section{What ``in \texorpdfstring{\(\iPRA\)}{iPRA}'' means}\label{sec:meaning}

\subsection{The arithmetic base}

The current ProveIt \code{ConciseIPRA} development supplies the syntax of
primitive-recursive functions \code{PR}, arithmetic terms \code{PTm}, the
axiom predicate \code{PRAAx}, intuitionistic derivability \code{Prf}, and the
quantifier-free induction constructor \code{PRA.qfInd}.  Its native
bi-interpretation exposes \code{NFm.ofPRA} and \code{pra_iff_native}.  The
object theory has only natural numbers.  Signed integers, rationals, finite
vectors, and analytic names below are codes.

A rational is a signed numerator and a positive denominator, reduced only when
convenient.  Equality and order are decided by cross multiplication.  A
Gaussian rational is a pair.  Finite lists are coded by a standard pairing
function; every fold has a displayed numerical bound.  Division is made total
by returning zero on a zero denominator, while every theorem that uses a
quotient also proves the denominator guard.

\subsection{Rapid dyadic names}

\begin{definition}[Regular rapid name]\label{def:rapid}
A real name is a fixed primitive-recursive dyadic sequence
\(p\mapsto x[p]\) with an integer bound \(B_x\ge1\) and a derivation of
\[
 \abs{x[p+1]-x[p]}\le 2^{-p-3},\qquad \abs{x[p]}\le B_x.
 \tag{2.1}
\]
A complex name is a pair of real names.  We write \(x=_\R y\) when
\[
 \forall p\;\abs{x[p+4]-y[p+4]}\le2^{-p}.
 \tag{2.2}
\]
An apartness witness \(x\apart0\) is finite data certifying that a rational
interval containing all later approximants lies wholly above \(2^{-m}\) or
wholly below \(-2^{-m}\), for some \(m\).
\end{definition}

Finite summation of (2.1) gives, for \(q\ge p\),
\[
 \abs{x[q]-x[p]}<2^{-p-2}.
 \tag{2.3}
\]
Thus the interval \(I_x(p)=[x[p]-2^{-p-2},x[p]+2^{-p-2}]\) contains every
later approximation.  Addition, multiplication on bounded names, reciprocal
from apartness, nonnegative square root, and rational powers are computed by
outward interval arithmetic.  Their regularity proofs are finite estimates.
For example,
\[
 \abs{xy-x'y'}\le \abs{x}\abs{y-y'}+\abs{y'}\abs{x-x'},
 \qquad
 \abs{x^{-1}-y^{-1}}\le \frac{2\abs{x-y}}{d^2}
 \tag{2.4}
\]
when \(d\le\abs{x},\abs{y}\).

\subsection{Exponential and logarithm with explicit tails}

For \(\abs z\le B\), put
\[
 E_N(z)=\sum_{j<N}\frac{z^j}{j!},\qquad
 N(B,p)=2B^2+2B+p+4.
\]
For \(j\ge2B\), successive majorant terms have ratio at most \(1/2\), and
\[
 \sum_{j\ge N(B,p)}\frac{B^j}{j!}<2^{-p-3}.
 \tag{2.5}
\]
The proof replaces the infinite tail by a finite geometric majorant; all
factorials, powers, and sums are primitive recursive.  Outward evaluation of
\(E_N\) at a rapid name defines \(\exp z\).  The finite binomial identity
\[
 \sum_{r+s=k}\frac{x^r y^s}{r!s!}=\frac{(x+y)^k}{k!}
\]
together with (2.5) proves \(e^{x+y}=e^x e^y\), positivity on real names, and
\(e^{-x}=1/e^x\).

For a positive name with an apartness lower bound, choose an integer \(k\) by
bounded comparison so that \(y=2^{-k}x\in[1,2]\).  With
\(u=(y-1)/(y+1)\), \(0\le u\le1/3\), define
\[
 \log y=2\sum_{j=0}^{M-1}\frac{u^{2j+1}}{2j+1}+R_M,
 \qquad
 0\le R_M\le\frac{2u^{2M+1}}{(2M+1)(1-u^2)}.
 \tag{2.6}
\]
Then \(\log x=k\log2+\log y\).  The tail bound supplies a primitive-recursive
modulus, and finite Cauchy products prove \(\log(ab)=\log a+\log b\).  The
inequalities used later,
\[
 \frac{t}{1+t}\le\log(1+t)\le t\quad(t>-1,\ t\ge0\text{ for the right bound}),
 \tag{2.7}
\]
and
\[
 \log(1-r)\le-r,
 \qquad
 \log(1-r)\ge-\frac{r}{1-r}\quad(0\le r<1),
 \tag{2.8}
\]
follow by differentiating finite Taylor polynomials and bounding their tails,
or directly by rational Riemann sums for \(1/(1+t)\).

\subsection{Derivatives and integrals as arithmetic statements}

A derivative assertion is an explicit modulus statement for rational difference
quotients.  For analytic names it is obtained by termwise differentiation plus
a majorant.  On a compact rational interval, an integral is the regular name
of midpoint Riemann sums; a supplied modulus of continuity determines the mesh.
An improper integral additionally carries two primitive-recursive tail
moduli.  Consequently substitution and the fundamental theorem are families of
finite error inequalities, not appeals to a completed continuum.

For later use, for positive rational \(a\) define
\[
 \Gamma(a)=\int_0^\infty t^{a-1}e^{-t}\,dt.
 \tag{2.9}
\]
The contribution of \([0,\delta]\) is at most \(\delta^a/a\); after a rational
cutoff, repeated integration by parts and (2.5) bound the tail.  This gives a
rapid name.  The finite integration-by-parts proof yields
\(\Gamma(a+1)=a\Gamma(a)\).  The usual constructive Gaussian-integral
argument supplies \(\Gamma(1/2)=\sqrt\pi\) after \(\pi\) is defined by the
located zero of cosine; the companion trigonometric packet contains that
construction.

\section{The map \texorpdfstring{\(F(w)=w e^w\)}{F(w)=w exp(w)} and the two real branches}\label{sec:realbranches}

Put \(F(w)=we^w\).  The finite product rule and the exponential derivative give
\[
 F'(w)=e^w(1+w),\qquad F''(w)=e^w(2+w).
 \tag{3.1}
\]
These are statements about difference quotients with explicit exponential
remainders.  Since \(e^w>0\), the sign of \(F'\) is the sign of \(1+w\).

\begin{lemma}[Constructive monotonicity and minimum]\label{lem:Fmono}
The function \(F\) is strictly decreasing on \(( -\infty,-1]\), strictly
increasing on \([-1,\infty)\), and has the unique global minimum
\[
 F(-1)=-e^{-1}.
\]
Moreover \(F(w)\to0^-\) as \(w\to-\infty\) and \(F(w)\to+\infty\) as
\(w\to+\infty\), with primitive-recursive rates.
\end{lemma}

\begin{proof}
For rational \(u<v\) in either interval,
\[
 F(v)-F(u)=\int_u^v e^t(1+t)\,dt.
\]
The integrand has the required sign; if the endpoints are positively separated
and the interval is not the singleton \(-1\), a rational subinterval has a
positive lower bound for its absolute value, proving strictness.  For
\(M\ge1\), \(0<M e^{-M}<2^{-p}\) once, for example,
\(M=4p+16\); use \(e>2\), proved from the first three terms of its series.
For \(w\ge0\), \(F(w)\ge w\).  These are explicit rates.
\end{proof}

Near the minimum a quadratic modulus replaces the missing inverse Lipschitz
constant.  If \(w=-1+s\), \(0\le s\le1\), then
\[
 eF(-1+s)+1=1+(s-1)e^s=\int_0^s t e^t\,dt,
\]
so
\[
 \frac{s^2}{2}\le eF(-1+s)+1\le\frac{e s^2}{2}.
 \tag{3.2}
\]
This elementary estimate is the key to a uniform inverse construction at the
branch point.

\begin{algorithm}[Certified inverse search]\label{alg:inverse-search}
For the principal branch and requested precision \(p\):
\begin{enumerate}
\item obtain a rational bracket \([-1,U]\), taking \(U=0\) for negative input
and \(U=B_x+1\) for nonnegative input;
\item subdivide the bracket into dyadic cells of width at most \(2^{-p-4}\);
\item evaluate \(F\) and the input by outward intervals at precision
\(2p+O(B_x+p)\), and retain the unique cell in which the monotone crossing can
occur; if two neighboring value intervals overlap, retain their union;
\item return the midpoint, with the retained cell as its error certificate.
\end{enumerate}
For the secondary branch, an input apartness witness
\(x\le-2^{-m}<0\) gives the left bracket \(-M\), with
\(M=4m+16\), since \(-M e^{-M}>-2^{-m}\).  Search the decreasing interval
\([-M,-1]\).
\end{algorithm}

All loops in Algorithm~\ref{alg:inverse-search} are bounded.  Estimate (3.2)
ensures that an unresolved comparison near \(-1/e\) still confines the inverse
to the requested width.  Thus equality at a grid point does not require an
undecidable real-number equality test.

\begin{theorem}[Real branches, domains, and ranges]\label{thm:branches}
The real solutions of \(we^w=x\) are exactly as follows.
\begin{enumerate}[label=\textup{(\roman*)}]
\item If \(x<-1/e\), there is no real solution.
\item If \(x=-1/e\), the unique solution is \(w=-1\).
\item If \(-1/e<x<0\), there are exactly two solutions:
\[
 \Wz(x)\in(-1,0),\qquad \Wm(x)\in(-\infty,-1).
\]
\item If \(x\ge0\), the unique solution is \(\Wz(x)\in[0,\infty)\).
\end{enumerate}
The map \(\Wz:[-1/e,\infty)\to[-1,\infty)\) is a strictly increasing
bijection.  The map \(\Wm:[-1/e,0)\to(-\infty,-1]\) is a strictly decreasing
bijection.  Both are primitive-recursively named operations on inputs carrying
the displayed domain certificates.
\end{theorem}

\begin{proof}
Lemma~\ref{lem:Fmono} gives injectivity and the endpoint values on each monotone
piece.  Algorithm~\ref{alg:inverse-search} supplies an inverse name and, by
continuity, its graph equation.  Any other solution in the same range equals it
by strict monotonicity.  The two ranges overlap only on \([-1/e,0)\), producing
exactly the stated multiplicities.
\end{proof}

\begin{corollary}[Continuity moduli]\label{cor:Wmoduli}
On every rational compact subinterval separated from \(-1/e\), each branch has
an explicit Lipschitz modulus obtained from a positive lower bound for
\(\abs{F'}\).  In an explicit rational neighbourhood of the common endpoint,
\[
 \abs{\Wz(x)+1}\le\sqrt{2e\abs{x+1/e}},
\]
and the analogous local estimate holds for \(\Wm\).  Thus both endpoint
extensions have a primitive-recursive square-root modulus.  The neighbourhood
restriction is essential for the secondary branch, since \(\Wm(x)\to-\infty\)
as \(x\uparrow0\).
\end{corollary}

\section{Identities and branch discipline}\label{sec:identities}

Every identity in this section is proved by substituting into \(we^w=x\) and
then invoking uniqueness on the displayed branch range.  This branch check is
not optional.

\begin{theorem}[Inverse identities]\label{thm:inverse-identities}
For real \(u\),
\[
 \Wz(ue^u)=u\quad(u\ge-1),
 \qquad
 \Wm(ue^u)=u\quad(u\le-1).
 \tag{4.1}
\]
For a branch value \(w=W(x)\), with the necessary nonzero guards,
\[
 e^w=\frac{x}{w},\qquad e^{-w}=\frac{w}{x},\qquad
 e^{nw}=\left(\frac{x}{w}\right)^n,
 \tag{4.2}
\]
and, for \(x>0\),
\[
 \log \Wz(x)=\log x-\Wz(x).
 \tag{4.3}
\]
\end{theorem}

\begin{proof}
Equation (4.1) follows because \(u\) lies in the selected inverse range.
Equation (4.2) is cancellation in \(we^w=x\); the guards produce the required
apartness witnesses.  Taking the constructive logarithm of the positive
identity \(we^w=x\) gives (4.3).
\end{proof}

\begin{theorem}[Logarithmic substitution identities]\label{thm:log-subst}
For positive \(x\),
\begin{align}
 \Wz(x\log x)&=\log x &&(x\ge1/e),\tag{4.4}\\
 \Wm(x\log x)&=\log x &&(0<x\le1/e),\tag{4.5}\\
 \Wz\!\left(-\frac{\log x}{x}\right)&=-\log x &&(0<x\le e),\tag{4.6}\\
 \Wm\!\left(-\frac{\log x}{x}\right)&=-\log x &&(x\ge e).\tag{4.7}
\end{align}
At \(x=e\), both displayed branch candidates meet at \(-1\).
\end{theorem}

\begin{proof}
Set \(u=\log x\).  Then \(ue^u=x\log x\).  The condition \(x\ge1/e\) is
exactly \(u\ge-1\), while \(0<x\le1/e\) is \(u\le-1\).  For the second pair,
put \(v=-\log x\); then \(ve^v=-\log x/x\), and the inequalities select the
proper range in (4.1).
\end{proof}

\begin{theorem}[Multiplication and addition identities]\label{thm:compound-identities}
For \(n\in\N_{>0}\) and \(x>0\),
\[
 \Wz\!\left(\frac{n x^n}{\Wz(x)^{n-1}}\right)=n\Wz(x).
 \tag{4.8}
\]
For \(x,y>0\),
\[
 \Wz(x)+\Wz(y)
 =\Wz\!\left(xy\left(\frac1{\Wz(x)}+\frac1{\Wz(y)}\right)\right).
 \tag{4.9}
\]
More general signed or complex versions hold after the branch is supplied as
part of the data.
\end{theorem}

\begin{proof}
Let \(a=\Wz(x)>0\).  Then the argument in (4.8) is
\(n(ae^a)^n/a^{n-1}=na e^{na}\), so uniqueness on the positive branch gives
the result.  For (4.9), writing \(a=\Wz(x)\), \(b=\Wz(y)\), the argument is
\[
 (ae^a)(be^b)(a^{-1}+b^{-1})=(a+b)e^{a+b}.
\]
Since \(a+b>0\), the principal branch returns \(a+b\).
\end{proof}

\begin{proposition}[Ratio parametrization of the two roots]\label{prop:ratio}
Let \(-1/e<c<0\), put \(a=\Wz(c)\), \(b=\Wm(c)\), and
\(t=b/a>1\).  Then
\[
 a=\frac{\log t}{1-t},\qquad
 b=\frac{t\log t}{1-t}.
 \tag{4.10}
\]
\end{proposition}

\begin{proof}
The equality \(ae^a=be^b\) and \(b=ta\) give
\(e^{(1-t)a}=t\).  Taking logarithms yields
\((1-t)a=\log t\), and multiplying by \(t\) gives the second formula.  Every
quantity has the displayed positive or negative separation, so the divisions
are constructive.
\end{proof}

\section{Special values and a branch-cut correction}\label{sec:special}

\begin{theorem}[Real special values]\label{thm:special-values}
The following identities hold with the stated branch choices:
\begin{align*}
 \Wz(-1/e)&=-1,& \Wz(0)&=0,& \Wz(e)&=1,\\
 \Wz(2\log2)&=\log2,&
 \Wz(e^{1+e})&=e,&
 \Wz(\sqrt e/2)&=1/2,\\
 \Wz(e^{1/n}/n)&=1/n\quad(n\ge1),&&
 \Wm(-\log2/2)&=-\log4.
\end{align*}
For every positive \(x\),
\[
 \Wz(x^{x+1}\log x)=x\log x.
 \tag{5.1}
\]
\end{theorem}

\begin{proof}
Each right-hand side \(u\) is substituted into \(ue^u\).  For example,
\((\log2)e^{\log2}=2\log2\),
\((1/n)e^{1/n}=e^{1/n}/n\), and
\((-\log4)e^{-\log4}=-(2\log2)/4=-\log2/2\).  The range condition selects the
branch.  For (5.1), set \(u=x\log x\); then
\(ue^u=x\log x\cdot x^x=x^{x+1}\log x\).  The elementary minimum
\(x\log x\ge-1/e>-1\) puts \(u\) in the principal range.
\end{proof}

\begin{definition}[Omega constant]
The rapid name \(\Omega\) is \(\Wz(1)\).  It is produced by the principal
inverse search on \([0,1]\).
\end{definition}

\begin{corollary}[Omega identities]\label{cor:omega}
\[
 \Omega e^\Omega=1,
 \qquad \Omega=e^{-\Omega},
 \qquad \Omega=-\log\Omega.
 \tag{5.2}
\]
A rational interval certificate in the accompanying verifier gives
\(0.5671432904097838<\Omega<0.5671432904097840\).
\end{corollary}

\begin{audit}[The value at \(-\pi/2\)]\label{audit:pi-cut}
Direct multiplication proves
\[
 \left(\frac{i\pi}{2}\right)e^{i\pi/2}=-\frac\pi2,
 \qquad
 \left(-\frac{i\pi}{2}\right)e^{-i\pi/2}=-\frac\pi2.
\]
Thus the two conjugate numbers are both Lambert roots.  Since
\(-\pi/2<-1/e\) lies on the standard principal cut, the notation
\(W_0(-\pi/2)\) is convention dependent there.  The upper boundary value is
\(i\pi/2\) and the lower boundary value is \(-i\pi/2\); the unqualified
single-valued equality on the reference page suppresses this choice.
\end{audit}

\section{A limited constructive complex-branch atlas}\label{sec:complex}

A global complex branch is not a primitive object of first-order arithmetic.
We instead fix a rational polygonal path \(\gamma\), a Gaussian-rational disk
containing a chosen root at its initial point, and a rational separation
\(\delta>0\) from \(w=-1\) along the lifted path.  Quantitative inverse
estimates subdivide \(\gamma\) into a primitive-recursively bounded number of
steps.  On each step, interval Newton produces the unique continued root.  This
constructs a named branch along that path and proves the defining equation and
derivative formula there.

\begin{proposition}[Preimage of the real axis]\label{prop:quadratrix}
For \(w=u+iv\),
\[
 \operatorname{Im}(we^w)=e^u(u\sin v+v\cos v).
\]
Away from \(\sin v=0\), the condition \(we^w\in\R\) is therefore
\[
 u=-v\cot v,
 \tag{5.3}
\]
and the corresponding real value is
\[
 we^w=-v\csc v\,e^{-v\cot v}.
 \tag{5.4}
\]
These are exact identities for constructive sine, cosine, and exponential
names whenever the displayed denominator has an apartness witness.
\end{proposition}

\begin{proof}
Expand
\((u+iv)e^u(\cos v+i\sin v)\), equate the imaginary part to zero, and
substitute (5.3) in the real part.  The calculation is finite Gaussian-rational
algebra after replacing the functions by their names.
\end{proof}

\begin{proposition}[Critical value]\label{prop:critical}
The only finite critical point of \(F(w)=we^w\) is \(w=-1\), and its critical
value is \(-1/e\).  Hence every named local branch separated from that value is
locally analytic, while the two real branches meet there with the square-root
modulus (3.2).
\end{proposition}

\begin{remark}[Conditional transcendence statement]
For a nonzero algebraic \(z\), if a value \(w\) of \(W(z)\) were nonzero and
algebraic, then \(e^w=z/w\) would be algebraic, contradicting
Lindemann--Weierstrass.  The reduction is elementary and formalizable in
\(\iPRA\).  A proof that the value is transcendental is labelled \StatusC here,
because this archive does not arithmetize an effective
Lindemann--Weierstrass theorem.
\end{remark}

\section{Differential calculus}\label{sec:calculus}

On a named branch separated from \(-1/e\), differentiate
\(W(z)e^{W(z)}=z\).  The finite product and chain rules give
\[
 e^W(1+W)W'=1.
 \]
Because \(z=We^W\), this has several equivalent forms.

\begin{theorem}[Differential equation and first derivative]\label{thm:derivative}
For every named branch at a point separated from the critical value,
\[
 z(1+W)W'=W,
 \tag{6.1}
\]
\[
 W'(z)=\frac{W(z)}{z(1+W(z))}
       =\frac1{z+e^{W(z)}}.
 \tag{6.2}
\]
At the removable point on the principal branch,
\[
 \Wz'(0)=1.
 \tag{6.3}
\]
\end{theorem}

\begin{proof}
The first two identities are the differentiated defining equation and
cancellation with explicit denominator guards.  At zero, write the defining
equation as \(W/z=e^{-W}\) for \(z\ne0\).  The branch continuity gives
\(W(z)\to0\), hence \(W(z)/z\to1\) with the exponential modulus.  This is the
difference quotient at zero.
\end{proof}

\begin{theorem}[All higher derivatives]\label{thm:higher-derivatives}
Define integer polynomials by
\[
 P_1(w)=1,
 \qquad
 P_{n+1}(w)=(1+w)P_n'(w)+(1-3n-nw)P_n(w).
 \tag{6.4}
\]
Then, for \(n\ge1\),
\[
 W^{(n)}(z)=
 \frac{P_n(W(z))}
 {(z+e^{W(z)})^n(1+W(z))^{n-1}}.
 \tag{6.5}
\]
The coefficient vector of \(P_n\) is a primitive-recursive function of \(n\).
Consequently the claim that a root \(w\) of \(P_n\) yields a root
\(we^w\) of \(W^{(n)}\) is a finite algebraic theorem whenever the denominator
is apart from zero.
\end{theorem}

\begin{proof}
Using \(z=we^w\), the denominator in (6.5) is
\(e^{nw}(1+w)^{2n-1}\).  Differentiation with respect to \(z\) is the operator
\[
 \frac{d}{dz}=\frac{e^{-w}}{1+w}\frac{d}{dw}.
\]
Applying it to
\(P_n(w)e^{-nw}(1+w)^{-(2n-1)}\) and collecting the numerator gives exactly
(6.4).  This is an induction on \(n\); polynomial differentiation and
multiplication are bounded vector operations.
\end{proof}

The first polynomials are
\[
 P_1=1,\qquad P_2=-(w+2),\qquad
 P_3=2w^2+8w+9,
\]
\[
 P_4=-(6w^3+36w^2+79w+64).
\]
The accompanying symbolic checker verifies the recurrence and (6.5) through a
user-selectable finite order.

\begin{corollary}[Second derivative and curvature]\label{cor:curvature}
For \(w=W(z)\),
\[
 W''(z)=-\frac{e^{-2w}(w+2)}{(w+1)^3}
       =-\frac{w^2(w+2)}{z^2(w+1)^3}.
 \tag{6.6}
\]
Hence \(W_0\) is strictly concave on \((-1/e,\infty)\).  On the secondary
branch, it is convex when \(-2<W_{-1}(x)<-1\), that is
\(-1/e<x<-2/e^2\), and concave when
\(W_{-1}(x)<-2\), that is \(-2/e^2<x<0\).  Its unique inflection point is
\((x,w)=(-2/e^2,-2)\).
\end{corollary}

\begin{proof}
Formula (6.6) is (6.5) for \(n=2\).  The sign calculation uses only the
ranges in Theorem~\ref{thm:branches}.  On the principal range, \(w+1>0\) and
\(w+2>0\), so the expression is negative.  On the secondary range the cubic
denominator is negative, so the sign equals that of \(w+2\).
\end{proof}

\section{Antiderivatives and definite integrals}\label{sec:integrals}

The substitution
\[
 w=W(x),\qquad x=we^w,
 \qquad dx=e^w(1+w)\,dw
 \tag{7.1}
\]
is constructive on every compact interval lying in one monotone branch.  It is
a theorem about Riemann-sum names with an explicit change-of-variables modulus.

\begin{theorem}[Elementary antiderivatives]\label{thm:antiderivatives}
On a fixed differentiable branch,
\begin{align}
 \int W(x)\,dx
 &=xW(x)-x+e^{W(x)}+C,
 \tag{7.2}\\
 \int \frac{W(x)}x\,dx
 &=\frac{W(x)^2}{2}+W(x)+C,
 \tag{7.3}\\
 \int W(Ae^{Bx})\,dx
 &=\frac{W(Ae^{Bx})^2}{2B}+\frac{W(Ae^{Bx})}{B}+C
 \quad(B\apart0),
 \tag{7.4}\\
 \int \frac{W(x)}{x^2}\,dx
 &=\Ei(-W(x))-e^{-W(x)}+C.
 \tag{7.5}
\end{align}
Here \(\Ei\) is the named antiderivative characterized on a separated interval
by \(\Ei'(t)=e^t/t\).
\end{theorem}

\begin{proof}
For (7.2), (7.1) changes the integrand to
\(w e^w(1+w)\); a direct derivative shows that
\(e^w(w^2-w+1)=xW-x+e^W\) is an antiderivative.  For (7.3),
\(W(x)dx/x=(1+w)dw\).  In (7.4), put \(w=W(Ae^{Bx})\).
Differentiating \(we^w=Ae^{Bx}\) gives
\(w'=Bw/(1+w)\), and the derivative of
\((w^2/2+w)/B\) is therefore \(w\).  Finally,
\[
 \frac{W(x)}{x^2}dx=e^{-w}\left(1+\frac1w\right)dw,
\]
and the derivative of \(\Ei(-w)-e^{-w}\) is the displayed integrand.  Each
calculation can alternatively be checked by differentiating the proposed
right-hand side, avoiding any global substitution theorem.
\end{proof}

\begin{corollary}[A basic definite integral]\label{cor:0e}
\[
 \int_0^e \Wz(x)\,dx=e-1.
 \tag{7.6}
\]
\end{corollary}

\begin{proof}
The antiderivative in (7.2) has value \(e\) at \(x=e\), since \(W_0(e)=1\),
and limiting value \(1\) at \(x=0\), since \(W_0(0)=0\).
\end{proof}

\begin{theorem}[Gamma family of improper integrals]\label{thm:gamma-integrals}
For every fixed integer \(N>1\),
\[
 \int_0^\infty \Wz(x^{-N})\,dx
 =N^{1-1/N}\Gamma(1-1/N).
 \tag{7.7}
\]
In particular,
\[
 \int_0^\infty \Wz(x^{-2})\,dx=\sqrt{2\pi}.
 \tag{7.8}
\]
All integrals in (7.7) have primitive-recursive tail moduli.
\end{theorem}

\begin{proof}
Put \(u=W_0(x^{-N})\).  Then
\(x=u^{-1/N}e^{-u/N}\), and reversing the endpoints gives
\[
 \int_0^\infty W_0(x^{-N})dx
 =\frac1N\int_0^\infty
 u^{-1/N}(1+u)e^{-u/N}\,du.
\]
Scale \(u=Nv\) and apply
\(\Gamma(a+1)=a\Gamma(a)\):
\begin{align*}
 I_N
 &=N^{-1/N}\Gamma(1-1/N)
   +N^{1-1/N}\Gamma(2-1/N)\\
 &=N^{1-1/N}\Gamma(1-1/N).
\end{align*}
The endpoint majorants are powers at zero and exponential tails at infinity,
so the substitution is justified by explicit finite cutoffs.  For \(N=2\),
use \(\Gamma(1/2)=\sqrt\pi\).
\end{proof}

\begin{corollary}[Two further displayed definite integrals]\label{cor:more-integrals}
\[
 \int_0^\infty\frac{\Wz(x)}{x\sqrt{x}}\,dx=2\sqrt{2\pi},
 \tag{7.9}
\]
\[
 \int_0^\pi \Wz(2\cot^2 t)\sec^2t\,dt=4\sqrt\pi.
 \tag{7.10}
\]
\end{corollary}

\begin{proof}
For (7.9), put \(u=W_0(x)\).  Then
\[
 \frac{W_0(x)}{x\sqrt{x}}dx
 =(u^{1/2}+u^{-1/2})e^{-u/2}\,du.
\]
Scaling by \(u=2v\) gives
\(2\sqrt2\Gamma(3/2)+\sqrt2\Gamma(1/2)=2\sqrt{2\pi}\).
For (7.10), split at \(\pi/2\), use \(s=\tan t\) on each half, and exploit
evenness:
\[
 2\int_0^\infty W_0(2/s^2)ds
 =2\sqrt2\int_0^\infty W_0(x^{-2})dx=4\sqrt\pi.
\]
Every singular endpoint is controlled by the same majorants used in
Theorem~\ref{thm:gamma-integrals}.
\end{proof}

\section{Finite Lagrange inversion and the Maclaurin series}\label{sec:series}

The companion ProveIt package develops Lagrange inversion entirely in finite
truncated rings.  We restate the exact portion needed here so that the Lambert
argument is self-contained.

For a commutative ring \(R\) and \(N\in\N\), let
\(T_N(R)=R[X]/(X^{N+1})\).  Its elements are vectors
\((a_0,\ldots,a_N)\); Cauchy multiplication, differentiation, powers, and
composition by a series with zero constant term are bounded sums.

\begin{lemma}[Finite fixed point]\label{lem:finite-fix}
For \(\Phi\in T_N(R)\), define
\[
 G_0=0,\qquad G_{j+1}=X\Phi(G_j).
\]
Then \(G_{N+1}\) is the unique \(G\in XT_N(R)\) satisfying
\(G=X\Phi(G)\).
\end{lemma}

\begin{proof}
If two series agree below degree \(m\), their images under
\(A\mapsto X\Phi(A)\) agree below degree \(m+1\).  Induction shows that the
\(j\)-th iterate has stabilized through degree \(j-1\).  After \(N+1\) steps
all represented coefficients are stable.  The same one-degree improvement
proves uniqueness coefficient by coefficient.  The invariant is a finite
conjunction of decidable rational equalities, hence quantifier free.
\end{proof}

\begin{theorem}[Finite Lagrange--Buermann formula]\label{thm:lagrange}
Let \(G=X\Phi(G)\) in \(T_N(R)\), and let \(H\) be a polynomial or truncated
series.  For \(1\le n\le N\), whenever division by \(n\) is available,
\[
 \coeff{n}{H(G)}=\frac1n\coeff{n-1}{H'(T)\Phi(T)^n}.
 \tag{8.1}
\]
In particular,
\[
 \coeff{n}{G}=\frac1n\coeff{n-1}{\Phi(T)^n}.
 \tag{8.2}
\]
\end{theorem}

\begin{proof}
Represent Laurent series by finite coefficient windows and define the residue
as the coefficient of \(T^{-1}\).  The finite identity
\[
 \operatorname{Res}_T A(B(T))B'(T)=\operatorname{Res}_U A(U)
\]
for \(B(T)=b_1T+\cdots\), \(b_1\) invertible, follows by checking monomials:
all terms except exponent \(-1\) are formal derivatives and have zero residue,
while \(B'/B\) has residue one.  Apply this to the inverse relation
\(X=G/\Phi(G)\), rewrite coefficient extraction as a residue, and use the
finite product rule.  No limiting series or contour occurs; every coefficient
window has a primitive-recursive size bound.  This is the residue proof in the
companion archive, and it translates directly to quantifier-free induction.
\end{proof}

For Lambert \(W\), the equation \(W=Xe^{-W}\) has
\(\Phi(T)=e^{-T}\), where exponential is truncated at degree \(N\).
Theorem~\ref{thm:lagrange} gives the coefficient in one line.

\begin{theorem}[Lambert series at zero]\label{thm:W-series}
For every \(N\), \(\iPRA\) proves in \(\Q[X]/(X^{N+1})\) that
\[
 W_0(X)\equiv
 \sum_{n=1}^{N}\frac{(-n)^{n-1}}{n!}X^n.
 \tag{8.3}
\]
Analytically, for \(\abs{x}<1/e\),
\[
 W_0(x)=\sum_{n\ge1}\frac{(-n)^{n-1}}{n!}x^n
 =x-x^2+\frac32x^3-\frac83x^4+\frac{125}{24}x^5-\cdots.
 \tag{8.4}
\]
\end{theorem}

\begin{proof}
Equation (8.2) gives
\[
 \coeff{n}{W}=\frac1n\coeff{n-1}{e^{-nT}}
 =\frac1n\frac{(-n)^{n-1}}{(n-1)!}
 =\frac{(-n)^{n-1}}{n!}.
\]
For fixed \(N\) this uses only the finitely many rational coefficients
\(1/j!\), \(j\le N\).  Analytic convergence and identification with the
principal inverse are proved in Theorem~\ref{thm:radius} below.
\end{proof}

\begin{theorem}[Integral and complex powers]\label{thm:powers}
For every fixed integer \(r\), as a Laurent series at zero,
\[
 W_0(x)^r
 =\sum_{n\ge r}
   \frac{-r(-n)^{n-r-1}}{(n-r)!}x^n.
 \tag{8.5}
\]
For a fixed constructive complex parameter \(r\),
\[
 \left(\frac{W_0(x)}x\right)^r=e^{-rW_0(x)}
 =1+\sum_{n\ge1}\frac{r(n+r)^{n-1}}{n!}(-x)^n,
 \qquad \abs{x}<1/e.
 \tag{8.6}
\]
In (8.5), the isolated expression \(0^0\) that occurs at \(n=0\) is read as
\(1\).  The \(n=0\) coefficient in the frequently printed unified notation
for (8.6) is likewise read as \(1\), including when \(r=0\).
\end{theorem}

\begin{proof}
Apply (8.1) with \(H(T)=T^r\) for positive \(r\):
\[
 \coeff{n}{W^r}
 =\frac{r}{n}\coeff{n-r}{e^{-nT}}
 =\frac{r}{n}\frac{(-n)^{n-r}}{(n-r)!}
 =\frac{-r(-n)^{n-r-1}}{(n-r)!}.
\]
Multiplication by the invertible unit \(W/x=e^{-W}\) extends the identity to
all integral Laurent powers.  For (8.6), apply (8.1) to
\(H(T)=e^{-rT}\), or expand the already established identity
\((W/x)^r=e^{-rW}\).  Coefficient extraction gives
\(r(n+r)^{n-1}(-1)^n/n!\).  The argument is polynomial in the parameter
\(r\) at each positive degree, so a fixed complex rapid name is admissible.
\end{proof}

For example,
\[
 W_0(x)^2=x^2-2x^3+4x^4-\frac{25}{3}x^5+18x^6-\cdots.
\]

\begin{theorem}[Exact centered radius and boundary convergence]\label{thm:radius}
The series (8.4) has radius of convergence exactly \(1/e\).  It converges
absolutely on the boundary circle \(\abs{x}=1/e\), and on the open disk it is
the principal analytic germ satisfying \(We^W=x\).
\end{theorem}

\begin{proof}
Let \(c_n=n^{n-1}/n!\).  Direct cancellation gives
\[
 \frac{c_{n+1}}{c_n}=\left(1+\frac1n\right)^{n-1}.
 \tag{8.7}
\]
The elementary binomial estimates used to define the rapid name of \(e\) give
an explicit modulus for convergence of (8.7) to \(e\).  The ratio test, stated
as finite upper and lower geometric comparisons, therefore gives radius
\(1/e\).

Inside the disk, finite multiplication of the series followed by an explicit
tail estimate proves \(We^W=x\).  The value is near zero and hence lies in the
principal local inverse range.  The radius cannot be larger: if the germ were
analytic through \(-1/e\), continuity would give \(W(-1/e)=-1\), and
differentiating \(F(W(z))=z\) there would give
\(F'(-1)W'(-1/e)=1\), i.e. \(0=1\).

Finally a finite Wallis/Stirling estimate supplies a rational constant \(C\)
with
\[
 c_ne^{-n}\le\frac{C}{(n+1)^{3/2}}.
 \tag{8.8}
\]
The proof expands the factorial quotient into products, bounds logarithms by
finite Taylor inequalities, and groups the remaining positive series in
dyadic blocks.  Since \(\sum(n+1)^{-3/2}\) has a primitive-recursive tail
modulus, (8.8) proves absolute boundary convergence.
\end{proof}

\section{Square-root expansion at the branch point}\label{sec:puiseux}

Put
\[
 p=\sqrt{2(1+ez)}.
 \tag{9.1}
\]
If \(w=-1+s\), then
\[
 p^2=2\bigl(1+(s-1)e^s\bigr)
 =s^2+\frac23s^3+\frac14s^4+\frac1{15}s^5+\cdots.
 \tag{9.2}
\]
This is a ramified formal inverse with a nonzero linear term after taking the
square root.  The finite ramified Lagrange construction computes any requested
coefficient by bounded rational arithmetic.

\begin{theorem}[Puiseux branches]\label{thm:puiseux}
For \(z\) near \(-1/e\), with the nonnegative real square root in (9.1),
\begin{align}
 W_0(z)&=-1+p-\frac{p^2}{3}+\frac{11p^3}{72}
       -\frac{43p^4}{540}+\frac{769p^5}{17280}+O(p^6),
 \tag{9.3}\\
 W_{-1}(z)&=-1-p-\frac{p^2}{3}-\frac{11p^3}{72}
       -\frac{43p^4}{540}-\frac{769p^5}{17280}+O(p^6).
 \tag{9.4}
\end{align}
For every fixed order, the \(O\)-term means an explicit rational constant and
primitive-recursive neighborhood on which the remainder is bounded by that
constant times the indicated power.
\end{theorem}

\begin{proof}
Seek \(s=\pm p+a_2p^2+a_3p^3+\cdots+a_Np^N\), substitute in (9.2), and
successively set the coefficients of \(p^3,\ldots,p^{N+1}\) to zero.  At each
stage the new coefficient occurs linearly with coefficient \(\pm2\), so the
recursion is primitive recursive and unique.  It gives the displayed numbers.
On a rational disk in the \(p\)-plane, the omitted exponential tail and the
finite substitution residual have explicit majorants.  Since the derivative
of the normalized ramified map is separated from zero, the quantitative
inverse lemma converts the residual into the stated remainder.  Replacing
\(p\) by \(-p\) changes exactly the odd coefficients and exchanges the two
real branches.
\end{proof}

\section{Logarithmic asymptotic expansions}\label{sec:asymptotic}

For \(x>1\), write
\[
 L_1=\log x,\qquad L_2=\log\log x.
\]
The equation \(we^w=x\) is equivalent to
\[
 w+\log w=L_1.
 \tag{10.1}
\]
Put \(w=L_1-L_2+s\), expand
\(\log w=L_2+\log(1-(L_2-s)/L_1)\), and revert the resulting formal series in
\(1/L_1\).  Let \(\genfrac{[}{]}{0pt}{}{n}{k}\) denote the unsigned Stirling
number of the first kind, generated by the bounded recurrence
\[
 \genfrac{[}{]}{0pt}{}{n+1}{k}
 =n\genfrac{[}{]}{0pt}{}{n}{k}+\genfrac{[}{]}{0pt}{}{n}{k-1}.
\]

\begin{theorem}[Every finite asymptotic order]\label{thm:asymptotic}
For each fixed \(N\ge1\), define
\[
 A_\ell(Y)=\sum_{m=1}^{\ell}
 \frac{(-1)^{\ell-m}}{m!}
 \genfrac{[}{]}{0pt}{}{\ell}{\ell-m+1}Y^m.
 \tag{10.2}
\]
There are primitive-recursive rational constants \(K_N,C_N\) such that, when
\(L_1\ge K_N\),
\[
 \left|W_0(x)-\left(L_1-L_2+
 \sum_{\ell=1}^{N}\frac{A_\ell(L_2)}{L_1^\ell}\right)\right|
 \le C_N\frac{(1+L_2)^{N+1}}{L_1^{N+1}}.
 \tag{10.3}
\]
The first terms are
\begin{align*}
 W_0(x)={}&L_1-L_2+\frac{L_2}{L_1}
 +\frac{L_2(-2+L_2)}{2L_1^2}
 +\frac{L_2(6-9L_2+2L_2^2)}{6L_1^3}\\
 &+\frac{L_2(-12+36L_2-22L_2^2+3L_2^3)}{12L_1^4}
 +O\!\left(\frac{(1+L_2)^5}{L_1^5}\right).
 \tag{10.4}
\end{align*}
In particular,
\(W_0(x)=\log x-\log\log x+o(1)\), in the finite witness sense supplied by
(10.3).
\end{theorem}

\begin{proof}
Formal substitution in (10.1) proves, by finite coefficient comparison, that
the expression in parentheses has residual
\(R_N\) bounded as on the right of (10.3).  The logarithm tail is bounded by a
geometric series once \((L_2+\abs{s})/L_1\le1/2\), enforced by \(K_N\).
The function \(H(w)=w+\log w\) has derivative
\(1+1/w\ge1\) for positive \(w\).  Therefore the difference between the exact
root and an approximate root is at most the absolute residual.  This is a
constructive residual-to-error estimate, and all coefficient and cutoff
functions are primitive recursive.
\end{proof}

\begin{theorem}[Secondary-branch asymptotics]\label{thm:wm-asymptotic}
Let \(x\to0^-\), put \(U=-\log(-x)>1\) and \(V=\log U\).  For every fixed
\(N\), \(-W_{-1}(x)\) has an expansion obtained from the same finite reversion
applied to
\[
 v-\log v=U,
\]
starting with
\[
 W_{-1}(x)=-U-V-\frac{V}{U}
 +\frac{V(-2+V)}{2U^2}
 +O\!\left(\frac{(1+V)^3}{U^3}\right).
 \tag{10.5}
\]
The remainder has an explicit primitive-recursive bound after a computable
threshold.
\end{theorem}

\begin{proof}
Write \(W_{-1}(x)=-v\).  The defining equation becomes
\(v e^{-v}=-x\), hence \(v-\log v=U\).  Substitute
\(v=U+V+s\), expand \(\log v\) as a finite Taylor polynomial with a geometric
remainder, and compare coefficients.  Since the derivative of
\(v-\log v\) is \(1-1/v\), it is at least \(1/2\) after \(v\ge2\), so twice
the residual bounds the root error.
\end{proof}

\section{Bounds for the principal branch}\label{sec:bounds-principal}

All bounds are proved by comparing a candidate \(a(x)\) with \(W_0(x)\)
through the increasing map \(F\):
\[
 a\le W_0(x)\iff a e^a\le x,
 \qquad
 W_0(x)\le a\iff x\le a e^a,
 \tag{11.1}
\]
provided \(a\ge-1\).  This turns every claim into elementary exponential and
logarithm inequalities with explicit finite Taylor certificates.

\begin{theorem}[Hoorfar--Hassani logarithmic bounds]\label{thm:HH}
For \(x\ge e\), put \(t=\log x\ge1\) and \(u=\log t\ge0\).  Then
\[
 t-u+\frac{u}{2t}
 \le W_0(x)
 \le t-u+\frac{e}{e-1}\frac{u}{t}.
 \tag{11.2}
\]
\end{theorem}

\begin{proof}
For the lower candidate
\(a=t-u+u/(2t)\), set \(q=u/t\).  Since
\(u/t\le1/e\) (the maximum of \(\log t/t\) occurs at \(t=e\)),
\(a>0\).  Write
\[
 \frac{a}{t}=1-r,
 \qquad r=q\left(1-\frac1{2t}\right)\ge\frac q2.
\]
Using \(\log(1-r)\le-r\),
\[
 \log a=u+\log(1-r)\le u-\frac{u}{2t}=t-a.
\]
Thus \(a+\log a\le t\), equivalently \(ae^a\le e^t=x\), proving the lower
bound.

For the upper candidate, set \(c=e/(e-1)\) and
\(b=t-u+cu/t\).  Now
\[
 \frac bt=1-r,
 \qquad r=q\left(1-\frac ct\right).
\]
If \(r\le0\), then \(\log(b/t)\ge0\ge-cq\).  If \(0<r<1\), use
\(\log(1-r)\ge-r/(1-r)\).  Since \(0<r\le q\le1/e\),
\[
 \frac{r}{1-r}\le\frac{q}{1-q}\le\frac{e}{e-1}q=cq.
\]
Hence in both cases
\(\log b=u+\log(b/t)\ge u-cu/t=t-b\).  Therefore
\(b+\log b\ge t\), so \(be^b\ge x\).  Every inequality used is (2.7) or
(2.8), with a rational tail certificate.
\end{proof}

The next pair, often presented as a more complicated formula, is simply a
one-step enclosure for a scalar correction.

\begin{theorem}[Iacono--Boyd bounds]\label{thm:IB}
For \(x\ge e\), let
\[
 q=\frac{\log\log x}{\log x},\qquad
 A=\log\left(\frac{x}{\log x}\right),
 \qquad h=-\log(1-q).
\]
Then
\[
 A+\frac{A}{1+A}h\le W_0(x)
 \le A+h-\log\left(1+\frac{h}{1+A}\right).
 \tag{11.3}
\]
Equivalently, these are exactly the two expanded expressions displayed on the
reference page.
\end{theorem}

\begin{proof}
Write \(W_0(x)=A+d\).  Since
\(A=\log x-\log\log x\) and \(A/\log x=1-q\), equation
\(w+\log w=\log x\) becomes
\[
 d+\log\left(1+\frac dA\right)=h.
 \tag{11.4}
\]
The left side is strictly increasing for \(d\ge0\).  Since
\(\log(1+s)\le s\), (11.4) gives
\(h\le d(1+1/A)\), hence
\[
 d\ge d_L:=\frac{A}{1+A}h.
\]
For the upper bound put
\(r=h/(1+A)\) and
\(d_U=h-\log(1+r)\).  The inequality \(\log(1+r)\le r\) gives
\(d_U\ge h-r=Ah/(1+A)=Ar\).  Therefore
\[
 \log(1+d_U/A)\ge\log(1+r).
\]
Therefore
\[
 d_U+\log(1+d_U/A)\ge h.
\]
Monotonicity in (11.4) yields \(d\le d_U\).
\end{proof}

\begin{theorem}[General one-parameter upper bound]\label{thm:general-upper}
For \(y>1/e\) and \(x\ge-1/e\),
\[
 W_0(x)\le
 \log\left(\frac{x+y}{1+\log y}\right),
 \tag{11.5}
\]
with equality exactly when \(x=y\log y\).  Taking \(y=x+1\) gives
\[
 W_0(x)\le
 \log\left(\frac{2x+1}{1+\log(x+1)}\right)
 \tag{11.6}
\]
whenever the displayed specialization is defined, in particular for
\(x\ge-1/e\).
\end{theorem}

\begin{proof}
Put
\[
 R=\frac{x+y}{1+\log y},\qquad s=\frac{R}{y},\qquad U=\log R.
\]
First \(U\ge-1\): at \(x=-1/e\) this is the inequality
\(ey-2-\log y\ge0\), which vanishes at \(y=1/e\) and has positive derivative
thereafter.  Now
\begin{align*}
 Ue^U-x
 &=R\log R-x\\
 &=y\bigl(s\log s-s+1\bigr).
\end{align*}
The function \(s\log s-s+1\) is nonnegative for \(s>0\), because its derivative
is \(\log s\) and its unique minimum is zero at \(s=1\).  Thus
\(F(U)\ge x\), and (11.1) proves (11.5).  Equality is equivalent to
\(s=1\), i.e. \(x=y\log y\).
\end{proof}

\begin{theorem}[Square- and cube-root enclosure near the branch point]\label{thm:root-bounds}
For \(-1/e\le x\le0\),
\[
 \sqrt{1+ex}-1\le W_0(x)\le\sqrt[3]{1+ex}-1.
 \tag{11.7}
\]
\end{theorem}

\begin{proof}
Write \(W_0(x)=t-1\), \(0\le t\le1\).  Then
\[
 q:=1+ex=1+(t-1)e^t.
 \tag{11.8}
\]
We prove \(t^3\le q\le t^2\).  The nonnegative power series
\[
 q=\sum_{n\ge2}\frac{n-1}{n!}t^n
 \tag{11.9}
\]
shows
\(q/t^2\le\sum_{n\ge2}(n-1)/n!=1\), hence \(q\le t^2\).
For the other inequality, let \(r(t)=q-t^3\).  Then
\[
 r'(t)=t(e^t-3t).
\]
The ratio \(e^t/t\) is strictly decreasing on \((0,1]\), so
\(e^t-3t\) changes sign at most once.  Since \(r(0)=r(1)=0\), \(r\) first
increases and then decreases, and is nonnegative throughout.  The sign-change
argument is implemented by a located rational bisection; alternatively a
Bernstein-polynomial certificate on \([0,1]\) is finite.  Taking monotone
square and cube roots gives (11.7).
\end{proof}

\section{Bounds for the secondary branch}\label{sec:bounds-secondary}

\begin{theorem}[Bounds at the common branch point]\label{thm:wm-near}
For \(u>0\),
\[
 -1-\sqrt{2u}-u
 <W_{-1}(-e^{-u-1})
 <-1-\sqrt{2u}-\frac23u.
 \tag{12.1}
\]
\end{theorem}

\begin{proof}
Put \(W_{-1}(-e^{-u-1})=-1-s\), \(s>0\).  The defining equation reduces to
\[
 u=s-\log(1+s).
 \tag{12.2}
\]
For the left inequality in (12.1), it suffices to prove
\(s<\sqrt{2u}+u\), or
\(\log(1+s)<\sqrt{2u}\).  Set \(v=\log(1+s)>0\).  Then
\(u=e^v-1-v\), and
\[
 2u-v^2=2\sum_{n\ge3}\frac{v^n}{n!}>0.
\]

For the right inequality it suffices to prove
\(s>\sqrt{2u}+2u/3\).  With the same \(v\), this becomes
\[
 \frac{e^v-1+2v}{3}>\sqrt{2(e^v-1-v)}.
 \tag{12.3}
\]
Both sides are positive.  After squaring and multiplying by nine, the
difference is
\[
 Q(v)=(e^v-1+2v)^2-18(e^v-1-v).
 \tag{12.4}
\]
Its Taylor expansion has zero coefficients through degree three and strictly
positive coefficients thereafter.  Indeed, direct expansion gives
\[
 Q(v)=\sum_{n\ge4}\frac{2^n+4n-20}{n!}v^n>0,
\tag{12.5a}
\]
because the numerator is positive for every $n\ge4$.  This coefficient
identity and the exponential tail majorant are finite and are also checked by
the accompanying exact verifier.  Thus (12.3) holds,
and translating back to \(-1-s\) gives (12.1).
\end{proof}

\begin{theorem}[Bounds as the secondary branch tends to zero]\label{thm:wm-far}
For \(u>1\),
\[
 -u-\log u-\frac{\log u}{u-1}
 \le W_{-1}(-e^{-u})
 \le -u-\log u-\frac{\log u}{u+\log u}.
 \tag{12.5}
\]
\end{theorem}

\begin{proof}
Write \(W_{-1}(-e^{-u})=-v\), so \(v>1\) and
\[
 v-\log v=u.
 \tag{12.6}
\]
Put \(v=u+\log u+d\), \(d>0\).  Equation (12.6) becomes
\[
 d=\log\left(1+\frac{\log u+d}{u}\right).
 \tag{12.7}
\]
Dropping \(d\) inside the increasing logarithm and using
\(\log(1+a)\ge a/(1+a)\) gives
\[
 d\ge\log\left(1+\frac{\log u}{u}\right)
 \ge\frac{\log u}{u+\log u}.
 \tag{12.8}
\]
Using \(\log(1+a)\le a\) in (12.7) gives
\[
 d\le\frac{\log u+d}{u},
 \qquad d\le\frac{\log u}{u-1}.
 \tag{12.9}
\]
Substitute these bounds in \(-v=-u-\log u-d\); the reversal of order under
negation yields (12.5).
\end{proof}

\begin{remark}[Constructive strictness]
The strict inequalities in (12.1) use the positive exponential tail after its
first omitted term.  Thus an explicit lower rational witness can be extracted
from any positive lower bound for \(u\).  The weak inequalities in (12.5) are
valid without deciding equality of real names.
\end{remark}

\section{Nested representations and representations not reconstructed}\label{sec:representations}

The defining equation may be read as either
\[
 W=x e^{-W}
 \qquad\text{or}\qquad
 W=\log(x/W).
 \tag{13.1}
\]
Iterating these fixed-point maps produces the familiar nested exponential and
nested logarithm.  A constructive statement must include a contraction
certificate rather than treating the ellipsis as an object.

\begin{theorem}[Certified nested exponential]\label{thm:nested-exp}
Let \(x\) be a real input on the principal branch, and suppose a rational
interval \(I\) and rational \(q<1\) are supplied such that
\(W_0(x)\in I\), \(g(I)\subseteq I\), and
\(\abs{g'(t)}\le q\) on \(I\), where \(g(t)=xe^{-t}\).  Then the finite
recursion
\[
 a_0\in I,\qquad a_{n+1}=\frac{x}{e^{a_n}}
 \tag{13.2}
\]
satisfies
\[
 \abs{a_n-W_0(x)}\le q^n\abs{a_0-W_0(x)}.
 \tag{13.3}
\]
Consequently its limit is the nested exponential
\[
 W_0(x)=\cfrac{x}{\exp\!\cfrac{x}{\exp\!\cfrac{x}{\ddots}}}.
 \tag{13.4}
\]
For every real \(x\) with \(\abs{W_0(x)}<1\) and a positive witness for that
strict inequality, the inverse-search enclosure can compute a finite entry
stage and an interval \(I\) satisfying the hypotheses.
\end{theorem}

\begin{proof}
The derivative is \(g'(t)=-xe^{-t}=-g(t)\).  The interval hypotheses and the
constructive mean-value estimate give
\(\abs{g(a)-g(b)}\le q\abs{a-b}\), so induction proves (13.3).  If only
\(\abs{W_0(x)}\le1-\delta\) is initially known, continuity of \(g'\) gives an
explicit neighborhood of the fixed point on which \(\abs{g'}\le1-\delta/2\).
The certified inverse name and finite iteration are searched until the current
interval enters that neighborhood.  For real inputs in the stated stability
range, the standard alternating monotone enclosure supplies that entry; all
searches are bounded by the requested output precision and \(\delta\).
\end{proof}

\begin{theorem}[Certified nested logarithm]\label{thm:nested-log}
Let \(x>0\), and suppose a positive interval \(I\) contains \(W_0(x)\), is
invariant under \(h(t)=\log(x/t)\), and has a rational lower bound
\(\abs t\ge d>1\).  Then
\[
 b_{n+1}=\log\frac{x}{b_n}
\]
converges to \(W_0(x)\) with error at most \(d^{-n}\) times the initial error.
Thus, with an explicit stability witness for \(\abs{W_0(x)}>1\),
\[
 W_0(x)=\log\cfrac{x}{\log\cfrac{x}{\log\cfrac{x}{\ddots}}}.
 \tag{13.5}
\]
\end{theorem}

\begin{proof}
On \(I\), \(\abs{h'(t)}=1/\abs t\le1/d<1\).  The same finite contraction
induction as in Theorem~\ref{thm:nested-exp} applies.  The inverse-search name
constructs a stable invariant interval from the positive separation of
\(\abs{W_0(x)}\) from one.
\end{proof}

\begin{audit}[Poisson, Kalugin--Jeffrey--Corless, and continued-fraction forms]
The reference page also displays two nontrivial integral representations and a
continued fraction with specialized coefficients.  Their finite truncations
are legitimate rational or quadrature objects, and the verifier can compare
initial Taylor coefficients.  This archive does not claim the global integral
identities or convergence of that continued fraction: a complete proof would
require a dedicated constructive contour/Poisson-kernel development or an
explicit continued-fraction recurrence and convergence majorant not given on
the page.  They are labelled \StatusD rather than inferred from numerical
agreement.
\end{audit}

\section{Solving transcendental equations}\label{sec:equations}

\begin{theorem}[Affine function plus exponential]\label{thm:affine-exp}
Let \(b,c\apart0\).  Every branch solution of
\[
 x=a+b e^{cx}
\]
is
\[
 x=a-\frac1c W_k(-bc e^{ac}),
 \tag{14.1}
\]
and conversely every defined right-hand side solves the equation.  If
\(a,c\apart0\), every branch solution of
\[
 ax+b=e^{cx}
\]
is
\[
 x=-\frac ba-\frac1c
 W_k\!\left(-\frac ca e^{-bc/a}\right).
 \tag{14.2}
\]
The real number of solutions is determined by whether the Lambert argument is
below, at, or above \(-1/e\), together with the selected branch ranges.
\end{theorem}

\begin{proof}
For the first equation set \(w=c(a-x)\).  Then
\[
 we^w=-bc e^{ac},
\]
which is equivalent to (14.1).  For the second set
\(y=x+b/a\).  The equation becomes
\(ay=e^{cy-bc/a}\), and multiplication by \(-c e^{-cy}/a\) gives
\((-cy)e^{-cy}=-(c/a)e^{-bc/a}\).  Solving for \(x\) gives (14.2).  Reversing
the algebra proves the converse; every division carries the displayed
apartness guard.
\end{proof}

\begin{example}[The equation \(2x+2=3^x\)]
Taking logarithms only through the positive base, or applying (14.2), gives the
two real solutions
\[
 x=-1-\frac{W_0(-\log3/6)}{\log3},
 \qquad
 x=-1-\frac{W_{-1}(-\log3/6)}{\log3}.
 \tag{14.3}
\]
The accompanying interval evaluator certifies respectively
\(-0.79012<x<-0.79010\) and \(1.44455<x<1.44457\).
\end{example}

\begin{theorem}[Inverse of \(x^x\) with all real branches]\label{thm:superroot}
For \(x>0\), the equation \(x^x=y\) is equivalent, after setting
\(u=\log x\), to
\[
 ue^u=\log y.
\]
Hence every positive solution is
\[
 x=e^{W_k(\log y)}
   =\frac{\log y}{W_k(\log y)},
 \tag{14.4}
\]
with the quotient interpreted by continuity at \(y=1\).  There are no positive
solutions for \(0<y<e^{-1/e}\), one at the minimum, two for
\(e^{-1/e}<y<1\), and one for \(y\ge1\).  For \(y>1\) only the principal
branch occurs.
\end{theorem}

\begin{proof}
Taking the constructive logarithm gives \(x\log x=\log y\).  With
\(u=\log x\), \(x=e^u\), so this is \(ue^u=\log y\).  The branch count follows
from Theorem~\ref{thm:branches}.  The minimum of \(x\log x\) is \(-1/e\) at
\(x=e^{-1}\), and exponentiation gives the threshold \(e^{-1/e}\).
\end{proof}

\begin{theorem}[The standard generalized equation]\label{thm:general-eqn}
For \(a_0,c\apart0\), every branch solution of
\[
 e^{-cx}=a_0(x-r)
\]
is
\[
 x=r+\frac1c W_k\!\left(\frac{c e^{-cr}}{a_0}\right).
 \tag{14.5}
\]
\end{theorem}

\begin{proof}
Put \(u=c(x-r)\).  Multiplying the equation by
\(c e^{cx}/a_0\) gives \(ue^u=ce^{-cr}/a_0\).  Solve for \(u\) and reverse
the substitutions.
\end{proof}

\begin{theorem}[Thermodynamic-equilibrium algebra]\label{thm:thermo}
Suppose positive \(T,K\) satisfy
\[
 \log K=\frac aT+b+c\log T,
 \qquad c\apart0.
\]
Then every branch solution is
\[
 T=\exp\!\left(
 W_k\!\left(-\frac ac e^{(b-\log K)/c}\right)
 +\frac{\log K-b}{c}\right).
 \tag{14.6}
\]
The branch classification of the Lambert argument gives zero, one, or two
positive algebraic candidates.  Whether the premise is the correct physical
equilibrium law is external to this theorem.
\end{theorem}

\begin{proof}
Set \(L=\log T\), multiply the original equation by \(T=e^L\), and rearrange:
\[
 -\frac ac e^{(b-\log K)/c}
 =\left(L+\frac{b-\log K}{c}\right)
 e^{L+(b-\log K)/c}.
\]
Apply a branch of \(W\) and exponentiate.
\end{proof}

\begin{theorem}[Wien-type stationary equation]\label{thm:wien}
For fixed integer \(D>0\), the equation
\[
 e^{-x}=1-\frac{x}{D}
\]
has the Lambert form
\[
 x=D+W_k(-De^{-D}).
 \tag{14.7}
\]
The branch yielding \(x=0\) is the algebraic endpoint solution; the other
admissible real branch gives the positive stationary point when it exists.
\end{theorem}

\begin{proof}
Set \(u=x-D\).  Then \(1-x/D=-u/D\), and the equation gives
\((-u)e^{u}=De^{-D}\), equivalently \(u e^u=-De^{-D}\).
\end{proof}

\begin{theorem}[Bernoulli--Todd equation]\label{thm:todd}
For
\[
 Y=\frac{X}{1-e^X},
 \tag{14.8}
\]
put \(V=Y-X\).  Then
\[
 Ve^V=Ye^Y.
 \tag{14.9}
\]
Excluding the removable limiting case \(Y=-1\), the nontrivial real solution is
\[
 X=\begin{cases}
 Y-W_0(Ye^Y)=W_{-1}(Ye^Y)-W_0(Ye^Y),&Y<-1,\\
 Y-W_{-1}(Ye^Y)=W_0(Ye^Y)-W_{-1}(Ye^Y),&-1<Y<0.
 \end{cases}
 \tag{14.10}
\]
\end{theorem}

\begin{proof}
Equation (14.8) is \(Y-X=Ye^X\).  With \(V=Y-X\) and
\(X=Y-V\), this becomes \(V=Ye^Y e^{-V}\), hence (14.9).  If \(Y<-1\), then
\(Y=W_{-1}(Ye^Y)\), so the other root is \(W_0(Ye^Y)\).  If
\(-1<Y<0\), the roles reverse.  Substitution gives (14.10).
\end{proof}

\begin{theorem}[Orthogonal-trajectory formula]\label{thm:orthogonal}
Every nonzero real solution of
\[
 \left(\frac1y+y\right)dy=
 \left(\frac1x-x\right)dx
 \tag{14.11}
\]
on a sign-fixed interval satisfies, for a constant \(C\),
\[
 y^2=W_0\!\left(x^2e^{-2C-x^2}\right),
 \tag{14.12}
\]
after absorbing the sign convention into \(C\).  Conversely, every branch-safe
square root of (14.12) solves (14.11).
\end{theorem}

\begin{proof}
Integration gives
\(\log\abs y+y^2/2=\log\abs x-x^2/2-C\).  Doubling and exponentiating yields
\[
 y^2 e^{y^2}=x^2e^{-2C-x^2}.
\]
The right side is nonnegative, so the unique nonnegative inverse is \(W_0\).
Differentiate the identity to prove the converse.
\end{proof}

\section{Rooted trees and finite combinatorics}\label{sec:trees}

Let \(t_n=n^{n-1}\) for \(n\ge1\), the number of rooted labeled trees on
\(n\) vertices.  A completely finite proof uses the Prufer code: deleting the
least labeled leaf at each step bijects labeled trees on \(n\) vertices with
words of length \(n-2\) over \(\{1,\ldots,n\}\).  Choosing the root multiplies
the count by \(n\).  Encoding and decoding are bounded loops, so the bijection
and count are \StatusF theorems.

Define the truncated exponential generating series
\[
 T_N(X)=\sum_{n=1}^{N}\frac{n^{n-1}}{n!}X^n.
\]

\begin{theorem}[Tree functional equation and Lambert series]\label{thm:trees}
For every \(N\),
\[
 T_N(X)\equiv X e^{T_N(X)}\pmod{X^{N+1}}.
 \tag{15.1}
\]
Consequently
\[
 W_0(X)=-T(-X)
 \tag{15.2}
\]
as a formal series, and the Lambert coefficient formula is equivalent to the
rooted Cayley count \(n^{n-1}\).
\end{theorem}

\begin{proof}
Combinatorially, a rooted labeled tree is a root together with a finite set of
rooted labeled subtrees.  Expanding the finite set construction gives the
coefficient of \(Xe^T\); equivalently, one can verify (15.1) directly using
the multinomial theorem and the Prufer count.  Replacing \(X\) by \(-X\) gives
\((-T(-X))e^{-T(-X)}=X\).  Uniqueness in
Lemma~\ref{lem:finite-fix} identifies this series with \(W_0\).  Comparing the
coefficient of \(X^n\) yields
\((-1)^{n-1}n^{n-1}/n!\), exactly (8.3).
\end{proof}

\section{Certified numerical evaluation}\label{sec:numerics}

Algorithm~\ref{alg:inverse-search} already gives a globally reliable method on
the two real branches.  Faster iterations are useful after a branch-safe
starting enclosure has been obtained.

\begin{theorem}[Newton and Halley formulae]\label{thm:newton-halley}
For \(f(w)=we^w-x\), Newton's iteration is
\[
 w_{j+1}=w_j-
 \frac{w_je^{w_j}-x}{e^{w_j}(1+w_j)}.
 \tag{16.1}
\]
Halley's iteration is
\[
 w_{j+1}=w_j-
 \frac{w_je^{w_j}-x}
 {e^{w_j}(1+w_j)-
 \dfrac{(w_j+2)(w_je^{w_j}-x)}{2(w_j+1)}}.
 \tag{16.2}
\]
Both are primitive-recursive operations on dyadic intervals whenever every
displayed denominator has a positive separation certificate.
\end{theorem}

\begin{proof}
Newton's formula is \(w-f(w)/f'(w)\), using
\(f'(w)=e^w(1+w)\).  Halley's formula
\[
 w-\frac{2ff'}{2(f')^2-ff''}
\]
with \(f''(w)=e^w(w+2)\) simplifies to (16.2).  Outward interval evaluation
of the exponential and rational operations is primitive recursive.
\end{proof}

\begin{theorem}[Logarithmic Newton iteration]\label{thm:log-newton}
On either real branch, when \(x/w>0\), \(w\apart0\), and \(w+1\apart0\), Newton's method for
\[
 g(w)=w+\log\abs w-\log\abs x
\]
is
\[
 w_{j+1}=
 \frac{w_j}{1+w_j}
 \left(1+\log\frac{x}{w_j}\right).
 \tag{16.3}
\]
Its fixed points are exactly the branch-safe roots of \(we^w=x\).
\end{theorem}

\begin{proof}
Since \(x/w_j>0\), the logarithm is defined.  Newton's update for
\(g(w)=w+\log\abs w-\log\abs x\), with
\(g'(w)=(1+w)/w\), is
\begin{align*}
 w-\frac{w}{1+w}
 (w+\log\abs w-\log\abs x)
 &=\frac{w}{1+w}\left(1+\log\frac{x}{w}\right).
\end{align*}
A fixed point with \(w\ne0\) gives
\(w+\log\abs w=\log\abs x\), and the common sign of \(x,w\) restores
\(we^w=x\).
\end{proof}

\begin{theorem}[Quadratic error certificate]\label{thm:newton-error}
Let a rational interval \(I\) contain a unique root \(\alpha\) of a twice
constructively differentiable function \(f\).  Suppose
\[
 \abs{f'(t)}\ge m>0,
 \qquad \abs{f''(t)}\le M
 \quad(t\in I),
 \tag{16.4}
\]
and the Newton update remains in \(I\).  Then
\[
 \abs{w_{j+1}-\alpha}
 \le\frac{M}{2m}\abs{w_j-\alpha}^2.
 \tag{16.5}
\]
For Lambert \(W\), rational bounds for \(m,M\) are obtained from
\(e^w(1+w)\) and \(e^w(2+w)\).  Thus the number of iterations required for
\(p\) output bits is computed in advance.
\end{theorem}

\begin{proof}
Taylor's formula with integral remainder, represented by finite Riemann sums,
gives
\[
 0=f(\alpha)=f(w)+f'(w)(\alpha-w)+R,
 \qquad \abs R\le\frac M2\abs{\alpha-w}^2.
\]
Subtract the Newton update and divide by the lower bound \(m\).  All estimates
are rational interval inequalities.
\end{proof}

\begin{algorithm}[Interval Newton with branch preservation]\label{alg:interval-newton}
Start from a rational branch bracket \(I=[a,b]\) supplied by
Algorithm~\ref{alg:inverse-search}.  Let \(m_I\) be an interval enclosure of
\(F'(I)\), separated from zero.  With midpoint \(c\), set
\[
 N(I)=c-\frac{F(c)-x}{m_I}
\]
using interval division, and replace \(I\) by \(I\cap N(I)\).  If the
intersection does not contract enough, take one bisection step.  Every update
preserves the root; away from the branch point the width becomes quadratic.
Near \(-1/e\), switch to the Puiseux coordinate \(p\) from (9.1) or retain
bisection.
\end{algorithm}

\subsection{Published starting values and the audit line}

The reference page lists branch-specific starting values for (16.3) and
explicit doubly exponential error estimates, including constants
\(\log(1+1/e)\), \((1-1/e)^{2^n-1}/5\), \(10^{-2^n}\), and
\(2^{-2^n}\).  Each such estimate is, in principle, a \StatusA theorem: after
substitution of the starting formula it reduces to finitely many inequalities
for \(\exp,\log,\sqrt{\cdot}\) on rational intervals.  This archive records the
formulas in \code{THEOREM_LEDGER.md} and numerically stress-tests them, but does
not claim a complete replay of all source-specific global constants.  The
generic certified scheme above supplies the same computational conclusion with
constants calculated from the actual input enclosure.

Likewise, the page's piecewise minimax rational approximants are finite
polynomial evaluations, but their quoted bit accuracy depends on external
coefficient tables and a uniform approximation proof.  They are diagnostics,
not native theorems of this package.

\section{Arithmetization in intuitionistic PRA}\label{sec:arithmetization}

\subsection{Primitive-recursive data and bounded loops}

The following operations are primitive recursive on natural-number codes:
\begin{enumerate}
\item signed integer, rational, dyadic, and Gaussian-rational arithmetic;
\item bounded vectors, Cauchy products, polynomial derivative, composition,
powers, and coefficient extraction;
\item factorials, binomial and Stirling numbers, and the polynomials \(P_n\);
\item outward interval evaluation of finite exponential and logarithm sums;
\item dyadic grid search, bisection, interval intersection, and a fixed number
of Newton or Halley steps;
\item finite Riemann sums and the explicit mesh/tail functions appearing in the
named integrals;
\item the Prufer encoder/decoder and finite tree-counting bijections.
\end{enumerate}

Each theorem is first proved at a requested precision or truncation degree.
The loop invariant is a quantifier-free formula about coded integers and
bounded lists.  Quantifier-free induction proves it for the loop bound.
Universal introduction over the precision parameter then yields extensional
equality of rapid names.

\subsection{No hidden classical principles}

\begin{center}
\begin{tabularx}{\textwidth}{>{\raggedright\arraybackslash}p{0.24\textwidth} X}
\toprule
Potential hidden principle & Constructive replacement\\
\midrule
Excluded middle for real order & Branch inputs carry located interval or
apartness certificates.  Bisection uses interval overlap and never asks whether
an arbitrary real is exactly zero.\\
Countable or dependent choice & Every approximation is a value of one fixed
primitive-recursive program; no sequence is selected from existential data.\\
Unbounded minimization & Every search has a bound computed from the magnitude,
apartness, precision, or truncation data.\\
Completed infinite series & A finite partial sum plus a primitive-recursive
tail bound defines each output approximation.\\
Classical inverse-function theorem & Monotone dyadic search and quantitative
residual-to-error estimates construct the inverse directly.\\
Division by a merely nonzero real & Every quotient uses positive apartness,
which computes a denominator lower bound.\\
Complex branch selection & A base root, a path, a cut convention, and positive
separation from the critical point are explicit inputs.\\
Improper integrals & Compact quadrature plus two explicit endpoint/tail moduli.\\
\bottomrule
\end{tabularx}
\end{center}

\begin{remark}[Schema discipline]
There is no primitive-recursive universal evaluator for all primitive-recursive
programs.  Therefore a theorem about a real name is a schema: the programs for
the inputs and their certificates are fixed before the first-order formula is
formed.  Uniformity remains available in ordinary numerical parameters such as
precision, branch flag, truncation degree, and finite coefficient vectors.
\end{remark}

\subsection{Suggested object-language formula}

At the lowest level, a principal-branch approximation theorem has the shape
\[
 \forall p\;\operatorname{W0Good}(xCode,p,w0Code(xCode,p)),
 \tag{17.1}
\]
where \code{W0Good} is quantifier free after bounded quantifiers are expanded:
it asserts a width bound, the range \([-1,U]\), and that the outward interval
for \(w e^w-x\) contains zero.  Uniqueness is a separate formula saying that
any two good intervals of sufficiently small width overlap after a fixed
precision shift.  The branch-point estimate (3.2) is the only non-Lipschitz
modulus needed.

For the coefficient theorem, the target is even more direct.  With a code for
truncated vectors,
\[
 \forall N\;\operatorname{mulTrunc}
 (W_N,\operatorname{expTrunc}(W_N,N),N)=X_N,
 \tag{17.2}
\]
where
\(W_N[n]=(-n)^{n-1}/n!\) for \(1\le n\le N\).  Clearing the common factorial
denominator turns (17.2) into integer equality, and the proof is quantifier-free
induction on the coefficient index.

\section{Lean formalization blueprint}\label{sec:lean}

The separate \code{lean/} directory contains a host-level analytic scaffold and
an object-language scaffold.  They deliberately contain \code{sorry}; no
kernel-checked completion is claimed.  The current environment used to build
this archive had no Lean executable, so syntax was reviewed but not compiled.

\subsection{Repository interfaces to reuse}

The target import is
\code{ConciseIPRA.NativeBiInterpretation}.  The relevant existing declarations
are \code{PR}, \code{PTm}, \code{PRAAx}, \code{Prf},
\code{PRA.qfInd}, \code{NFm.ofPRA}, and \code{pra_iff_native}.  The prior
Lagrange-inversion packet should be factored into reusable modules rather than
reproved inside Lambert \(W\).

\subsection{Recommended module order}

\begin{longtable}{>{\ttfamily}p{0.30\textwidth} p{0.62\textwidth}}
\toprule
Module & Main declarations and proof obligations\\
\midrule
\endhead
Codes.lean & Signed integers, rationals, dyadics, Gaussian rationals, bounded
vectors, interval arithmetic, total division plus guards.\\
RapidName.lean & Regular names, extensional equality, apartness, arithmetic,
square root, compact interval maps.\\
ExpLog.lean & Truncated exponential/logarithm evaluators, tail bounds,
addition/multiplication laws, derivatives, inequalities (2.7)--(2.8).\\
FiniteSeries.lean & Truncated Cauchy product, composition, derivative, formal
residue, finite Lagrange--Buermann.  Reuse the existing packet.\\
LambertReal.lean & \code{F}, monotonicity, branch-point quadratic modulus,
certified inverse search, \code{W0Name}, \code{Wm1Name}, uniqueness and ranges.\\
LambertSeries.lean & Coefficients, powers, radius majorants, Puiseux recursion,
asymptotic coefficient generator.\\
LambertCalculus.lean & First and higher derivatives, polynomial recurrence,
curvature, antiderivative checks.\\
LambertBounds.lean & Logarithm inequalities and the candidate comparisons in
Sections~11--12.\\
LambertIntegrals.lean & Effective Riemann sums, gamma name, change of variables,
formulas (7.6)--(7.10).\\
LambertNumerics.lean & Bisection, interval Newton, Newton/Halley/log-Newton,
residual-to-error theorem.\\
LambertApplications.lean & Equation transformations, Prufer codes, rooted tree
series, Bernoulli--Todd and orthogonal trajectories.\\
LambertObject.lean & \code{PR} codes, quantifier-free correctness formulas,
\code{PRA.qfInd} proofs, and transfer through \code{pra_iff_native}.\\
\bottomrule
\end{longtable}

\subsection{Representative host-level interfaces}

\begin{lstlisting}[language={}] 
structure RapidName where
  approx : Nat -> Rat
  bound  : Nat
  regular : forall p,
    |approx (p+1) - approx p| <= 1 / (2:Rat)^(p+3)
  bounded : forall p, |approx p| <= bound

def ExtEq (x y : RapidName) : Prop :=
  forall p, |x.approx (p+4) - y.approx (p+4)| <= 1/(2:Rat)^p

structure Apart (x : RapidName) where
  precision : Nat
  exponent  : Nat
  sign      : Bool
  separated : IntervalSeparated x precision exponent sign

structure BranchSpec (x w : RapidName) where
  graph : ExtEq (mul w (expName w)) x
  range : PrincipalRange w          -- or SecondaryRange w

theorem W0_unique (h1 : BranchSpec x u) (h2 : BranchSpec x v) :
  ExtEq u v := by
  -- monotonicity plus the quadratic endpoint modulus
  sorry
\end{lstlisting}

\subsection{Representative finite-series target}

\begin{lstlisting}[language={}]
def lambertCoeff (n : Nat) : Rat :=
  if h : n = 0 then 0
  else ((-(n : Int))^(n-1) : Rat) / factorial n

theorem lambert_trunc_graph (N : Nat) :
  mulTrunc (lambertVec N)
    (expTrunc (lambertVec N) N) N = Xvec N := by
  -- finite Lagrange inversion, or direct coefficient induction
  sorry
\end{lstlisting}

After a native theorem is proved, the object-language file should instantiate
the corresponding \code{PR} term, prove its defining equations with
\code{PRA.qfInd}, and only then use \code{pra_iff_native}.  Treating a host
Mathlib theorem as if it were already an object-theory derivation would be a
category error.

\section{Coverage audit against the reference page}\label{sec:audit}

Table~\ref{tab:audit} is deliberately conservative.  ``Proved'' means that the
paper contains the human argument and an explicit route to a first-order
\(\iPRA\) derivation.  The machine-readable ledger in the archive separates
individual formulas more finely.

\begin{longtable}{>{\raggedright\arraybackslash}p{0.34\textwidth}
                  >{\centering\arraybackslash}p{0.08\textwidth}
                  >{\raggedright\arraybackslash}p{0.49\textwidth}}
\caption{Coverage of the mathematical assertions on the Lambert \(W\) page.}
\label{tab:audit}\\
\toprule
Claim family & Status & Result or qualification\\
\midrule
\endfirsthead
\toprule
Claim family & Status & Result or qualification\\
\midrule
\endhead
Definition as inverse of \(we^w\) & \StatusA & Constructed by certified
monotone interval search; Theorem~\ref{thm:branches}.\\
Two real branches, domains, ranges, multiplicities & \StatusA & Complete proof,
including endpoint moduli; Theorem~\ref{thm:branches}.\\
Complex branch notation and cuts & \StatusA/\StatusD & Local path-based branches
with explicit separation are constructed; a global Riemann-surface
formalization is not claimed.\\
Critical point and branch point & \StatusA & Proposition~\ref{prop:critical}
and Puiseux Theorem~\ref{thm:puiseux}.\\
Image/preimage curve for real values & \StatusF/\StatusA & Quadratrix equations
(5.3)--(5.4), Proposition~\ref{prop:quadratrix}.\\
Transcendence at nonzero algebraic arguments & \StatusC & Reduced to effective
Lindemann--Weierstrass; that theorem is not developed here.\\
Differential equation and first derivative & \StatusA & Equations
(6.1)--(6.3), including the removable value at zero.\\
Higher derivative polynomial recurrence & \StatusF/\StatusA & Full induction,
Theorem~\ref{thm:higher-derivatives}.\\
Monotonicity and curvature & \StatusA & Lemma~\ref{lem:Fmono} and
Corollary~\ref{cor:curvature}.\\
Antiderivative of \(W\) and \(\int_0^eW_0\) & \StatusA & Equations
(7.2) and (7.6).\\
Other displayed indefinite integrals & \StatusA & Equations (7.3)--(7.5), with
guards and an effective \(\Ei\) name.\\
Maclaurin series of \(W_0\) & \StatusF/\StatusA & Finite coefficient theorem
and analytic lifting, Theorem~\ref{thm:W-series}.\\
Radius \(1/e\) and boundary convergence & \StatusA & Ratio modulus, critical
obstruction, and \(n^{-3/2}\) majorant; Theorem~\ref{thm:radius}.\\
Integral powers of \(W_0\) & \StatusF/\StatusA & Laurent coefficient formula
(8.5).\\
Complex powers of \(W_0/x\) & \StatusF/\StatusA & Formula (8.6) for a fixed
constructive complex parameter.\\
Large-\(x\) logarithmic asymptotic expansion & \StatusF/\StatusA & Every fixed
order with a residual bound, Theorem~\ref{thm:asymptotic}.\\
\(W_{-1}\) asymptotic as \(x\to0^-\) & \StatusF/\StatusA & Every fixed order;
Theorem~\ref{thm:wm-asymptotic}.\\
Hoorfar--Hassani bounds & \StatusA & Short proof in
Theorem~\ref{thm:HH}.\\
Iacono--Boyd enhanced bounds & \StatusA & Reduced to the correction equation
(11.4), Theorem~\ref{thm:IB}.\\
General one-parameter upper bound and specialization & \StatusA &
Theorem~\ref{thm:general-upper}.\\
Square/cube-root bounds on \([-1/e,0]\) & \StatusA &
Theorem~\ref{thm:root-bounds}.\\
Both secondary-branch bound families & \StatusA & Theorems
\ref{thm:wm-near} and \ref{thm:wm-far}.\\
Defining, logarithmic, multiplication, and addition identities & \StatusA &
Theorems~\ref{thm:inverse-identities}--\ref{thm:compound-identities}.\\
Branch-ratio parametrization & \StatusA & Proposition~\ref{prop:ratio}.\\
Real special values and omega identities & \StatusA &
Theorem~\ref{thm:special-values} and Corollary~\ref{cor:omega}.\\
\(W_0(-\pi/2)=i\pi/2\) & Audit & Correct only after an upper-boundary branch
convention; Audit~\ref{audit:pi-cut}.\\
Poisson integral representation & \StatusD & Not reconstructed; requires a
constructive Poisson/contour certificate.\\
Kalugin--Jeffrey--Corless integral representation & \StatusD & Not
reconstructed; finite quadrature tests are not a proof.\\
Displayed continued fraction & \StatusD & Initial convergents can be checked,
but no general recurrence/convergence certificate was supplied.\\
Nested exponential and logarithm & \StatusA & Proved with explicit invariant
interval and contraction witness; Theorems~\ref{thm:nested-exp} and
\ref{thm:nested-log}.\\
Gamma-family improper integrals & \StatusA & Theorem~\ref{thm:gamma-integrals}
and Corollary~\ref{cor:more-integrals}.\\
Affine/exponential equations & \StatusF/\StatusA & Exact branch-safe formulas,
Theorem~\ref{thm:affine-exp}.\\
Inverse of \(x^x\) & \StatusA & Corrected full branch count,
Theorem~\ref{thm:superroot}.\\
Rooted-tree EGF and Cayley coefficients & \StatusF/\StatusA & Prufer bijection
and finite series identity, Theorem~\ref{thm:trees}.\\
Bernoulli--Todd equation & \StatusF/\StatusA & Theorem~\ref{thm:todd}.\\
Thermodynamic formula & \StatusF/\StatusA & Algebraic solution of the displayed
equation, not the physical premise; Theorem~\ref{thm:thermo}.\\
Wien stationary equation & \StatusF/\StatusA & Theorem~\ref{thm:wien}.\\
Orthogonal trajectories & \StatusA & Integrated and checked by differentiation;
Theorem~\ref{thm:orthogonal}.\\
General equation \(e^{-cx}=a_0(x-r)\) & \StatusF/\StatusA &
Theorem~\ref{thm:general-eqn}.\\
Newton, Halley, logarithmic Newton & \StatusF/\StatusA & Formulae, guards, and
generic certified convergence; Section~\ref{sec:numerics}.\\
Published universal starting-value error constants & \StatusP & Recorded and
stress-tested, but the source-specific global inequality certificates are not
fully replayed.\\
Piecewise minimax rational approximants & \StatusD & Evaluation is finite; the
quoted uniform bit accuracy needs the external coefficient tables and proof.\\
Software availability table & \StatusD & A time-sensitive implementation list,
not a mathematical theorem.\\
Physical and engineering occurrence claims & \StatusD & The displayed algebraic
solutions may be checked, but the model derivations require their cited domain
literature.\\
Generalized Lambert functions with quadratic/infinite products & \StatusD & New
functions, outside the standard inverse treated here.\\
\bottomrule
\end{longtable}

\section{Validation, limitations, and reproducibility}\label{sec:validation}

The script \code{scripts/verify\_lambert\_w.py} performs exact rational checks
of the following finite assertions:
\begin{enumerate}
\item the first configurable number of coefficients of
\(W e^W=X\);
\item the power formulae (8.5)--(8.6) for several integer parameters;
\item the recurrence (6.4) and higher-derivative formula through a configurable
order;
\item the Puiseux coefficients in (9.3)--(9.4);
\item the positive coefficient identity (12.5a);
\item the first asymptotic polynomials from the Stirling formula;
\item representative real-branch identities, bounds, integrals, and iteration
residuals at high precision.
\end{enumerate}
Exact rational or symbolic checks are evidence for the finite algebra.  Floating
point tests are explicitly labelled diagnostics and are not substituted for a
proof.  The generated \code{evidence/verification\_report.txt} records the
parameters and results.  In the archived run, all 55 exact checks and 130
high-precision diagnostics passed (185/185 total).

The PDF is built from this source by \code{latexmk -pdf}.  The archive includes
a checksum ledger.  The rendered PDF was inspected page by page for clipping,
overlap, and broken glyphs.  No Lean executable was available in the build
environment; therefore the Lean blueprints were not compiled, and this fact is
recorded both here and in \code{lean/STATUS.md}.

\subsection{Remaining mathematical limitations}

The principal remaining formalization burden is not Lambert-specific algebra.
It is reusable infrastructure: an efficient object-language library for rapid
names and interval arithmetic, effective compact integration with change of
variables, and path-based complex analytic continuation.  Once those layers
exist, most theorems in this paper become short compositions of generic lemmas.
The highly specialized integral representations and continued fraction would
still require separate analytic developments.

\section{Further research questions}\label{sec:further}

\begin{enumerate}
\item Can the real-branch inverse be implemented in \code{ConciseIPRA} with a
single common interval-search combinator that also serves inverse trigonometric
and algebraic functions?
\item What is the smallest collection of quantitative analytic lemmas needed to
turn the local path-lifting construction into all integer-indexed complex
branches on standard slit domains?
\item Can the \(n^{-3/2}\) boundary majorant for the Lambert series be proved by
a shorter Wallis-product certificate than a general constructive Stirling
formula?
\item Is there a direct finite combinatorial proof of the power coefficients
(8.5)--(8.6), extending the rooted-tree interpretation beyond the first power?
\item Can the Hoorfar--Hassani and Iacono--Boyd inequalities be generated
automatically by a verified ``candidate through monotone inverse'' tactic?
\item Can all published starting-value constants for logarithmic Newton be
replayed by interval certificates small enough for a practical \(\iPRA\)
proof checker?
\item Which of the Poisson and Kalugin--Jeffrey--Corless integral
representations has the shortest constructive proof after the trigonometric and
Riemann-integration packets are available?
\item Can an effective Lindemann--Weierstrass theorem be isolated in a form
weak enough for \(\iPRA\) yet strong enough to prove transcendence of
\(W_k(\alpha)\) for nonzero algebraic \(\alpha\)?
\item Can the branch-point Puiseux evaluator and logarithmic asymptotic evaluator
be combined into a formally verified near-optimal real implementation with a
small certificate language?
\item Which scientific applications on the reference page reduce to a common
verified affine-exponential equation compiler, and which genuinely require a
generalized Lambert function?
\end{enumerate}

\section{Conclusion}

The Lambert \(W\) function is well suited to a constrained constructive
foundation.  Its global real behavior comes from two monotone pieces of one
explicit elementary map; its local series is a finite Lagrange-inversion
identity; its branch point is a finite ramified inversion; its asymptotics are
formal reversion plus a residual bound; and its numerical evaluation is bounded
interval arithmetic.  None of these steps requires excluded middle for real
numbers, countable choice, or an impredicative completed function space.

The resulting package proves most of the reference page's reusable mathematical
formulas in the precise named-data sense appropriate to \(\iPRA\), while
leaving a visible boundary around global complex topology, deep transcendence,
specialized representations, and external scientific claims.  That boundary is
not a defect: it identifies the exact reusable libraries needed for a complete
Lean and object-theory development.

\section*{Artifact inventory}
\addcontentsline{toc}{section}{Artifact inventory}
\begin{itemize}
\item \code{lambert\_w\_ipra.tex} and \code{lambert\_w\_ipra.pdf}: this report.
\item \code{THEOREM\_LEDGER.md} and \code{claims.json}: formula-by-formula status,
guards, and proof locations.
\item \code{lean/LambertWHostBlueprint.lean}: host-level data structures and
analytic theorem skeletons.
\item \code{lean/LambertWObjectBlueprint.lean}: intended integration with
\code{PR}, \code{PRAAx}, \code{PRA.qfInd}, and \code{pra\_iff\_native}.
\item \code{lean/STATUS.md}: exact formalization status and recommended build
order.
\item \code{scripts/verify\_lambert\_w.py}: exact symbolic checks and numerical
diagnostics.
\item \code{evidence/verification\_report.txt}: generated check results.
\item \code{README.md}, \code{SOURCE\_AUDIT.md}, and \code{SHA256SUMS}.
\end{itemize}

\begin{thebibliography}{99}
\bibitem{wikipedia}
Wikipedia contributors,
\emph{Lambert W function}, revision inspected 16 November 2026.
\url{https://en.wikipedia.org/wiki/Lambert_W_function}

\bibitem{corless}
R. M. Corless, G. H. Gonnet, D. E. G. Hare, D. J. Jeffrey, and D. E. Knuth,
``On the Lambert W function,''
\emph{Advances in Computational Mathematics} 5 (1996), 329--359.

\bibitem{hoorfar}
A. Hoorfar and M. Hassani,
``Inequalities on the Lambert W function and hyperpower function,''
\emph{J. Inequal. Pure Appl. Math.} 9 (2008), Article 51.

\bibitem{iacono}
R. Iacono and J. P. Boyd,
``New approximations to the principal real-valued branch of the Lambert
W-function,'' \emph{Advances in Computational Mathematics} 43 (2017),
1403--1436.

\bibitem{loczi}
L. Loczi,
``Guaranteed- and high-precision evaluation of the Lambert W function,''
\emph{Applied Mathematics and Computation} 433 (2022), 127406.

\bibitem{chatz}
I. Chatzigeorgiou,
``Bounds on the Lambert function and their application to the outage analysis
of user cooperation,'' \emph{IEEE Communications Letters} 17 (2013),
1505--1508.

\bibitem{proveit}
V. Reshetnikov,
\emph{ProveIt}, especially
\code{Logic/PrimitiveRecursiveArithmetic/ConciseIPRA} and the constructive
analysis proof packets, repository inspected 16 November 2026.
\url{https://github.com/VladimirReshetnikov/ProveIt}

\bibitem{lagrangePacket}
\emph{Lagrange inversion in intuitionistic primitive recursive arithmetic},
companion artifact in the user's ProveIt research corpus, 15 November 2026.

\bibitem{trigPacket}
\emph{Trigonometric identities in intuitionistic primitive recursive
arithmetic}, companion artifact supplying rapid names, exponential, sine,
cosine, and \(\pi\), 2026.
\end{thebibliography}

\end{document}
